\section{Dérivations et calculs complémentaires}
\label{annexes:details_calculs}

\subsection{Simplification des équations d'Euler}
Cette dérivation concerne le cas non visqueux : l'hypothèse
$\boldsymbol{\tau} = \mathbf{0}$ dans les bilans
\eqref{eq:navier_stokes} donne les équations d'Euler.

La masse volumique est supposée uniforme dans l'espace
et constante dans le temps :
\begin{equation}
    \rho(\mathbf{x}, t) = \rho_0 > 0.
\end{equation}
La gravité est également négligée : $\mathbf{g} = \mathbf{0}$.

\paragraph*{Conservation de la masse}

L'équation de conservation de la masse s'écrit :
\begin{equation}
    \frac{\partial \rho}{\partial t}
    + \nabla \cdot (\rho \mathbf{u}) = 0.
\end{equation}

Le développement de la divergence du produit donne :
\begin{equation}
    \frac{\partial \rho}{\partial t}
    + \mathbf{u} \cdot \nabla \rho
    + \rho \nabla \cdot \mathbf{u} = 0.
\end{equation}

La masse volumique étant constante, ses dérivées temporelle
et spatiales sont nulles :
\begin{equation}
    \frac{\partial \rho}{\partial t} = 0,
    \qquad
    \nabla \rho = \mathbf{0}.
\end{equation}

Il reste donc :
\begin{equation}
    \rho_0 \nabla \cdot \mathbf{u} = 0.
\end{equation}

Puisque $\rho_0 > 0$, il en résulte la contrainte
d'incompressibilité :
\begin{equation}
    \boxed{\nabla \cdot \mathbf{u} = 0.}
\end{equation}

\paragraph*{Conservation de la quantité de mouvement}

L'équation de conservation de la quantité de mouvement s'écrit :
\begin{equation}
    \frac{\partial (\rho \mathbf{u})}{\partial t}
    + \nabla \cdot \left(
        \rho \mathbf{u} \otimes \mathbf{u}
        + p \mathbf{I}
    \right)
    = \rho \mathbf{g}.
\end{equation}

La linéarité de la divergence et l'identité
$\nabla \cdot (p \mathbf{I}) = \nabla p$ donnent :
\begin{equation}
    \frac{\partial (\rho \mathbf{u})}{\partial t}
    + \nabla \cdot
        \left(\rho \mathbf{u} \otimes \mathbf{u}\right)
    + \nabla p
    = \rho \mathbf{g}.
\end{equation}

La masse volumique étant constante et la gravité étant négligée,
cette équation devient :
\begin{equation}
    \rho_0 \frac{\partial \mathbf{u}}{\partial t}
    + \rho_0 \nabla \cdot
        \left(\mathbf{u} \otimes \mathbf{u}\right)
    + \nabla p
    = \mathbf{0}.
\end{equation}

La division par $\rho_0$ et le passage du gradient de pression
dans le membre de droite donnent :
\begin{equation}
    \frac{\partial \mathbf{u}}{\partial t}
    + \nabla \cdot \left(\mathbf{u} \otimes \mathbf{u}\right)
    = -\frac{1}{\rho_0}\nabla p.
\end{equation}

\paragraph*{Système obtenu}

Les deux équations constituent
le système \eqref{eq:euler_simplified} :
\begin{equation*}
    \begin{aligned}
        \nabla \cdot \mathbf{u} &= 0, \\
        \frac{\partial \mathbf{u}}{\partial t}
        + \nabla \cdot \left(\mathbf{u} \otimes \mathbf{u}\right)
        &= -\frac{1}{\rho_0}\nabla p.
    \end{aligned}
\end{equation*}

